\section{A shallower teacher and many additional learner layers}
\label{sec:padding}
The generic formula resource already has polynomial learner-depth dependence when additional layers only relay existing computations. A fixed small L\'evy intensity gives a sharper bound: the cost of those extra layers is nearly linear.

In this section the activation is ReLU or tanh. Both are $1$-Lipschitz, vanish at zero, and satisfy $|\phi(t)|\le|t|$. Choose the same fixed prior as before with
\begin{equation}
 M_1=\int|u|\,\nu(du)=\frac{2c}{p-1}\le\frac14.
 \label{eq:padding-intensity}
\end{equation}
This is only a choice of the scalar intensity constant. It imposes no support restriction on a neural network.

Let a teacher have hidden depth $h\ge1$, width $m$, coefficient cap $B\ge1$, and actual rooted formula size at most $S$. A learner has any prescribed depth $H\ge h$ and width $N\ge m$. Put
\[
 q=H-h,\qquad S'=S+mq,\qquad D_0=C(m+1)(B+8).
\]
Define
\begin{align*}
 \mathcal J_n=C(d+1)(S'+1)\big[&(h+1)\log D_0\\
 &+\log\{C(S'+1)(d+1)(q+1)(B+8)(n+1)\}\big].
\end{align*}
Constants in this section are independent of $h,H,m,B,N,n$.

\begin{theorem}[Nearly linear overhead for extra learner depth]
\label{thm:stable-padding}
For either ReLU or tanh and the intensity choice~\eqref{eq:padding-intensity}, the ordinary final posterior mean at temperature $\beta_*=1/[256(1+\sigma^2)]$ satisfies
\[
 \E_{S_n}R(\bar f_{S_n,\beta_*})
 \le R(g)+\frac{C(1+\sigma^2)(1+\mathcal J_n)}{n+1}.
\]
The high-probability companion at $\beta_\eta=\eta/[32(1+\sigma^2)]$ satisfies
\[
 R(\bar f_{S_n,\beta_\eta})
 \le R_0+(1+\eta)\{R(g)-R_0\}
 +\frac{C(1+\sigma^2)}{\eta n}
       \left[\mathcal J_n+\log\frac2\delta\right]
\]
with probability at least $1-\delta$. Both statements hold uniformly over every prescribed finite $N\ge m$ and $H\ge h$. For an arbitrary width-$m$ teacher one may take $S=\sum_{j=1}^h m^j$. The only nearly linear dependence on the extra learner depth is through $S'=S+m(H-h)$ and a logarithm of $H-h+1$.
\end{theorem}

\begin{proof}
It suffices to construct, for each $0<t\le1$, a data-independent comparison with
\begin{align}
 &\E_Q\sup_x|f_w(x)-g(x)|\le t,\nonumber\\
 &\KL(Q\|\mu_N)\le C(d+1)(S'+1)
 \left[(h+1)\log D_0+
 \log\frac{C(S'+1)(d+1)(q+1)(B+8)}t\right].
 \label{eq:padding-comparison}
\end{align}
For $q=0$ this follows from the uniform comparison proof already given.
Assume $q\ge1$. Unroll the teacher and append $q$ single-predecessor relay layers to each relevant last-layer feature. The number of added occurrences is at most $mq$. Use the same disjoint embedding construction as before; the small-width branch uses the original teacher with relay layers. It remains to control the relay precision without an exponentially growing stability factor.

Fix an embedding. At the end of the teacher portion, the absolute path bounds and its expected uniform error obey
\[
 S_0:=\max_{\rm selected}\E R_j\le U,\qquad
 P_0:=N^{-1}\sum_j\E R_j\le U,\qquad
 e_0\le U\epsilon,\qquad
 U=C(d+1)D_0^{h+1}.
\]
These follow from the previous path and uniform-error lemmas, enlarging $D_0$ by a fixed constant.

For a relay coefficient equal to one, take $\epsilon\le1/8$.
The designated law then has positive support and mean one, hence
$\E|W|=1$. At an unselected native row, independence gives
$\E R_j\le M_1P_{\rm prev}$. A selected relay row therefore satisfies
\begin{equation}
 S_\ell\le S_{\ell-1}+C\epsilon P_{\ell-1},\qquad
 P_\ell\le M_1P_{\ell-1}+S_\ell.
 \label{eq:relay-first-moments}
\end{equation}
The same bounds hold for the first relay if its designated coefficient lies in $(0,1]$ and its entire support is positive. Since $M_1\le1/4$, induction gives
\[
 S_\ell\le U(1+C'\epsilon)^\ell,\qquad
 P_\ell\le2U(1+C'\epsilon)^\ell.
\]
In particular these are bounded by $C U$ for all $\ell\le q$ whenever
$\epsilon q\le1$. This uses only first moments, and no layerwise support condition.

For ReLU, use exact identity relays. Its teacher last-layer features are nonnegative, so the target is unchanged. Expected uniform error satisfies
\[
 e_\ell\le e_{\ell-1}
    +C\epsilon(S_{\ell-1}+P_{\ell-1}),
\]
by Lipschitzness, independence, and the first absolute moments of designated perturbations and background weights. The readout then gives
\[
 \E\sup_x|f_w(x)-g(x)|
 \le C L_\psi m(B+1)(q+1)U\epsilon.
\]
Choose $\epsilon$ to make this at most $t$, also enforcing
$\epsilon\le1/8$ and $\epsilon q\le1$. Its inverse logarithm is at most a fixed multiple of
\[
 (h+1)\log D_0+\log\frac{C(d+1)(q+1)(B+8)}t.
\]
The coefficient cap remains $B$. Substitution in the embedding entropy bound proves~\eqref{eq:padding-comparison}.

For tanh, let $u$ denote a teacher last-layer feature, so $|u|\le1$.
In the first extra layer use coefficient $\delta$ on its designated edge, in the remaining relay layers use one, and multiply the corresponding final readout coefficient by $\delta^{-1}$. Choose
\[
 \delta=\min\left\{1,\sqrt{\frac{t}{C L_\psi m(B+1)(q+1)}}\right\}.
\]
The ideal relay value is $\delta u$. Since
$|\tanh v-v|\le|v|^3/3$, $q$ deterministic relays have total raw error at most
$q\delta^3/3$ before the final readout amplification.
Take $\epsilon\le\delta/8$. The first relay's designated law has positive support and mean $\delta\le1$, so~\eqref{eq:relay-first-moments} continues to hold. Comparing each selected relay directly to the ideal value $\delta u$ gives
\[
 e_q\le\delta e_0+CqU\epsilon+q\delta^3/3.
\]
After the readout, including its own perturbations and background,
\[
 \E\sup_x|f_w(x)-g(x)|
 \le C L_\psi m(B+1)
       \{U\epsilon+qU\epsilon/\delta+q\delta^2\}.
\]
Choose the absolute constant in $\delta$ sufficiently large and then take
$\epsilon\le c\,t\delta/[L_\psi m(B+1)(q+1)U]$, also imposing the preceding cutoff requirements. The right side is at most $t$.
The designated coefficient cap is at most $B/\delta+1$. Both
$\log(1/\epsilon)$ and $\log(B/\delta+9)$ are bounded by a fixed multiple of the same logarithmic expression displayed above. The comparison entropy again gives~\eqref{eq:padding-comparison}. The learner's prior and likelihood do not depend on these proof-only relay coefficients or cutoffs.

For the expected-risk statement take $t=(n+1)^{-1}$.
Because outputs and $f_0$ lie in $[-1,1]$, the pointwise squared-risk difference averaged over the comparison is at most $4t$. Apply the generic sharp expectation theorem and~\eqref{eq:padding-comparison}.
For the high-probability statement take $t=1/n$ and apply the relative-loss PAC--Bayes and fixed-comparison concentration lemmas. Here the comparison excess is at most $A(g)+4/n$, so the same proof as Theorem~\ref{thm:main} gives its displayed remainder. This completes both claims.
\end{proof}
